\documentclass[12pt]{article}
\usepackage[utf8]{inputenc}
\usepackage[russian]{babel}
\usepackage{amsmath, amssymb}
\usepackage{geometry}
\usepackage{enumitem}
\usepackage{titlesec}
\usepackage{fancyhdr}
\usepackage{xcolor}
\usepackage{tcolorbox}
\usepackage{array}
\usepackage{tabularx}
\usepackage{tikz}
\usepackage[colorlinks=true, linkcolor=blue, urlcolor=blue, citecolor=blue]{hyperref}
\tcbuselibrary{skins}

\geometry{margin=2cm}
\titleformat{\section}{\Large\bfseries}{\thesection.}{0.5em}{}
\titleformat{\subsection}{\normalsize\bfseries}{\thesubsection.}{0.5em}{}

\pagestyle{fancy}
\fancyhf{}
\fancyhead[C]{Силлабус курса «Дифференциальные уравнения» — Университет ИТМО}
\fancyfoot[C]{\thepage}

\setlength{\parindent}{15pt}

% Зелёный блок
\definecolor{greenbordercolor}{HTML}{129f57}
\newtcolorbox{greenblock}[1][]{
    enhanced,
    colback=greenbordercolor!05,
    left=3mm,
    right=3mm,
    top=2mm,
    bottom=2mm,
    frame code={
        \draw[line width=1mm, color=greenbordercolor] (frame.south west) -- (frame.north west);
    },
    #1
}

\begin{document}

\begin{center}
    {\LARGE\textbf{Силлабус курса «Дифференциальные уравнения»}}\\[0.3em]
    {\large Двухсеместровый профильно-адаптированный курс для бакалавриата Университета ИТМО}\\[0.4em]
\end{center}

\begin{greenblock}
\textbf{Цель курса:} изучение методов решения дифференциальных уравнений, теорем существования и единственности, качественной теории устойчивости, теории хаоса и вариационного исчисления. Курс рассчитан на 2 семестра и ориентирован на формирование практических вычислительных навыков на Python с индивидуальной адаптацией лабораторных работ под образовательные программы бакалавриата Университета ИТМО.
\end{greenblock}

\begin{greenblock}
\textbf{Методические принципы и статус реализации под программы ИТМО:}
\begin{itemize}[nosep, leftmargin=1.2em]
    \item \textbf{Профилирование лабораторных работ:} контекст практических заданий адаптируется под специфику конкретного направления подготовки ИТМО без снижения фундаментального математического уровня.
    \item \textbf{Текущий статус внедрения (Критерий 1):} в полном объёме разработаны, апробированы и реализованы в виде интерактивных `Jupyter Notebooks` лабораторные комплексы по направлениям \textbf{«Мехатроника и робототехника»} и \textbf{«Информационная безопасность»}.
    \item \textbf{Потенциал масштабирования (Критерий 4):} для остальных направлений бакалавриата ИТМО сформированы профильные вариационные модули, запланированные к внедрению в 2026/2027 уч.~году.
    \item \textbf{Аудиторный формат практикума:} лабораторные работы выполняются непосредственно на семинарских занятиях под контролем преподавателя.
    \item \textbf{Python-верификация:} каждое задание сопровождается эталонным кодом (SciPy, SymPy, Matplotlib) для численно-графической самопроверки.
\end{itemize}
\end{greenblock}

\section*{Матрица адаптации лабораторного практикума по направлениям ИТМО}\label{sec:matrix}

\begin{greenblock}
Ниже представлена сводная матрица соответствия официальных направлений подготовки бакалавриата Университета ИТМО и профильных модулей лабораторных работ с явной разметкой статуса готовности (\textbf{Реализовано} / \textbf{Запланировано}). Подробные математические постановки приведены в \hyperref[app:A]{\textbf{Приложении А}}.
\end{greenblock}

\vspace{0.3em}
\begin{center}
\small
\begin{tabularx}{\textwidth}{|p{4.2cm}|X|X|c|}
\hline
\textbf{Направление подготовки бакалавриата ИТМО} & \textbf{Семестр 1: Лабораторные №1--2} & \textbf{Семестр 2: Лабораторные №3--4} & \textbf{Статус} \\
\hline
\textbf{Мехатроника и робототехника} & Расчет динамики звеньев манипулятора и линеаризация роторного маятника & Вариационная оптимизация траектории робота на сложной поверхности & \textcolor{greenbordercolor}{\textbf{Реализовано}} \\
\hline
\textbf{Информационная безопасность} & Численные схемы и анализ устойчивости линейных систем киберзащиты & Потоковое XOR-шифрование на основе хаоса Лоренца, SHA-256 и функции Ляпунова & \textcolor{greenbordercolor}{\textbf{Реализовано}} \\
\hline
\textbf{Прикладная математика и информатика} & Численные методы ОДУ, анализ жесткости и линеаризация нелинейных систем & Показатели Ляпунова хаоса Лоренца и геодезические линии на поверхностях & \textcolor{blue}{\textbf{Запланировано}} \\
\hline
\textbf{Программная инженерия} (профиль «Нейротехнологии») & Непрерывные нейронные сети Хопфилда и фазовый портрет нейрона ФитцХью--Нагумо & Обучение Neural ODE и вариационный Adjoint-метод оптимизации весов & \textcolor{blue}{\textbf{Запланировано}} \\
\hline
\textbf{Прикладная информатика} (профиль «Мобильные технологии») & Динамика сетевого трафика и анализ фазовых портретов сетевых буферов & Моделирование детерминированного хаоса задержек пакетов TCP и оптимизация маршрутизации & \textcolor{blue}{\textbf{Запланировано}} \\
\hline
\textbf{Бизнес-информатика} & Численный анализ модели Басса (диффузия инноваций) и модели Гудвина & Вариационная оптимизация производственного выпуска и уровня товарных запасов & \textcolor{blue}{\textbf{Запланировано}} \\
\hline
\end{tabularx}
\end{center}

\newpage

\begin{center}
{\Large\textbf{ОСНОВНАЯ ПРОГРАММА КУРСА}}
\end{center}

\section*{СЕМЕСТР 1. Аналитические и численные методы ОДУ}

\subsection*{Раздел 1. Дифференциальные уравнения первого порядка}
\begin{itemize}[nosep]
    \item \textbf{§1. Введение:} Постановка задач. Классификация уравнений. Геометрическая интерпретация.
    \item \textbf{§2. Формы записи и задача Коши:} Уравнение в нормальной и симметрической форме. Понятие решения. Метод ломаных Эйлера.
    \item \textbf{§3. Уравнения, решаемые в квадратурах:} Разделяющиеся, однородные, линейные уравнения. Уравнения Бернулли и Риккати. Полные дифференциалы и интегрирующий множитель.
    \item \textbf{§4. Особые типы уравнений:} Уравнения, не разрешённые относительно производной. Метод введения параметра. Уравнения Лагранжа и Клеро.
    \item \textbf{§5. Теория существования и единственности:} Приближения Пикара, метод сжимающих отображений. Теоремы существования и единственности. Особые решения.
\end{itemize}

\subsection*{Раздел 2. Уравнения высших порядков и системы}
\begin{itemize}[nosep]
    \item \textbf{§6. Уравнения $n$-го порядка:} Нормальная форма. Постановка задачи Коши. Методы понижения порядка.
    \item \textbf{§7. Линейные уравнения и системы:} Построение фундаментальной системы решений. Уравнения с постоянными коэффициентами. Принцип суперпозиции. Краевые задачи и теорема Штурма.
    \item \textbf{§8. Качественный анализ и линеаризация:} Положения равновесия автономных систем. Линеаризация в окрестности точки покоя. Классификация особых точек (узел, фокус, седло, центр).
\end{itemize}

\begin{greenblock}
\textbf{Краткое содержание лабораторного практикума 1 семестра:}
\begin{itemize}[nosep]
    \item \textbf{Лабораторная работа №1:} Сравнительный анализ методов Эйлера и Рунге--Кутты 2-го порядка на профильных моделях (см. \hyperref[app:A1]{\textbf{Приложение А.1}}).
    \item \textbf{Лабораторная работа №2:} Линеаризация нелинейных систем и качественный анализ фазовых портретов на профильных моделях (см. \hyperref[app:A2]{\textbf{Приложение А.2}}).
\end{itemize}
\end{greenblock}

\section*{СЕМЕСТР 2. Устойчивость, хаос и вариационное исчисление}

\subsection*{Раздел 3. Теория устойчивости и элементы теории хаоса}
\begin{itemize}[nosep]
    \item \textbf{§9. Теория устойчивости Ляпунова:} Устойчивость по Ляпунову. Прямой метод Ляпунова. Асимптотическая устойчивость. Предельные циклы и бифуркации.
    \item \textbf{§10. Нелинейная динамика и детерминированный хаос:} Чувствительность к начальным условиям. Аттрактор Лоренца. Системы реакция--диффузия.
\end{itemize}

\subsection*{Раздел 4. Вариационное исчисление и оптимизация}
\begin{itemize}[nosep]
    \item \textbf{§11. Основы вариационного исчисления:} Функционал и его экстремум. Уравнение Эйлера--Лагранжа. Дискретизация функционалов и численная оптимизация (градиентный спуск, метод Ньютона).
\end{itemize}

\begin{greenblock}
\textbf{Краткое содержание лабораторного практикума 2 семестра:}
\begin{itemize}[nosep]
    \item \textbf{Лабораторная работа №3:} Хаотические системы, функции Ляпунова и профильные приложения (см. \hyperref[app:A3]{\textbf{Приложение А.3}}).
    \item \textbf{Лабораторная работа №4:} Дискретизация вариационных задач и профильная оптимизация (см. \hyperref[app:A4]{\textbf{Приложение А.4}}).
\end{itemize}
\end{greenblock}

\newpage

\begin{center}
{\Large\textbf{ПРИЛОЖЕНИЯ: ПОДРОБНЫЕ ПРОФИЛЬНЫЕ МОДУЛИ И СИСТЕМА ОЦЕНИВАНИЯ}}
\end{center}

\section*{Приложение А. Подробное описание профильных примеров лабораторных работ}\label{app:A}

\subsection*{А.1. Лабораторная работа №1. Численные методы ОДУ первого порядка}\label{app:A1}
\textit{Базовое задание:} Реализация методов Эйлера $\psi(t_{k+1}) = \psi(t_k) + h_k f(t_k, \psi(t_k))$ и Рунге--Кутты 2-го порядка для задачи Коши $x'(t) = f(t, x(t)), x(t_0) = x_0$.
\begin{itemize}[nosep]
    \item \textbf{Мехатроника и робототехника} \textcolor{greenbordercolor}{\textbf{[Реализовано]}}: Численный расчет переходных процессов скорости электропривода колесного робота.
    \item \textbf{Информационная безопасность} \textcolor{greenbordercolor}{\textbf{[Реализовано]}}: Численное решение уравнений динамического распределения ресурсов средств киберзащиты.
    \item \textbf{Прикладная математика и информатика} \textcolor{blue}{\textbf{[Запланировано]}}: Анализ жестких уравнений и оценка критериев численной устойчивости шага $h_k$.
    \item \textbf{Бизнес-информатика} \textcolor{blue}{\textbf{[Запланировано]}}: Моделирование распространения продуктов по модели диффузии Басса $\frac{dN(t)}{dt} = \left(p + q \frac{N(t)}{M}\right)(M - N(t))$ и капитала компании $\frac{dK}{dt} = I(t) - \delta K$. Численный поиск момента максимального роста продаж $\max N'(t)$.
    \item \textbf{Программная инженерия (Нейротехнологии)} \textcolor{blue}{\textbf{[Запланировано]}}: Непрерывная нейронная сеть Хопфилда $\dot{x}_i = -\beta x_i + \text{softmax}(Ax)_i$ для очистки цифровых изображений от шума и распознавания по тени.
    \item \textbf{Прикладная информатика (Мобильные технологии)} \textcolor{blue}{\textbf{[Запланировано]}}: Моделирование динамики очереди пакетов в буфере мобильного маршрутизатора.
\end{itemize}

\subsection*{А.2. Лабораторная работа №2. Линеаризация и фазовые портреты}\label{app:A2}
\textit{Базовое задание:} Вычисление матрицы Якоби $J(x^*)$, линеаризация нелинейных систем в окрестности точек покоя, построение фазовых портретов.
\begin{itemize}[nosep]
    \item \textbf{Мехатроника и робототехника} \textcolor{greenbordercolor}{\textbf{[Реализовано]}}: Динамика роторного маятника $\ddot{\theta} + \gamma\dot{\theta} + \omega_0^2\sin\theta = 0$. Анализ точек покоя $(0, 0)$ (устойчивый фокус) и $(\pi, 0)$ (седло), построение фазового портрета в плоскости $(\theta, \dot{\theta})$.
    \item \textbf{Информационная безопасность} \textcolor{greenbordercolor}{\textbf{[Реализовано]}}: Фазовый портрет нелинейного хаотического генератора Чуа (Chua's circuit) $\dot{x} = \alpha(y - x - g(x)), \; \dot{y} = x - y + z, \; \dot{z} = -\beta y$ и анализ возникновения бифуркаций.
    \item \textbf{Бизнес-информатика} \textcolor{blue}{\textbf{[Запланировано]}}: Модель рыночной конкуренции Гудвина / Лотки--Вольтерры $\begin{cases} \dot{x} = x(\alpha - \beta y) \\ \dot{y} = -y(\gamma - \delta x) \end{cases}$. Линеаризация в окрестности экономических колебаний и определение типа точки покоя (центр).
    \item \textbf{Программная инженерия (Нейротехнологии)} \textcolor{blue}{\textbf{[Запланировано]}}: Модель импульсного нейрона ФитцХью--Нагумо $\dot{v} = v - v^3/3 - w + I, \; \dot{w} = \epsilon(v + a - bw)$. Построение фазовой плоскости $(v, w)$, анализ кубической изоклины и автоколебательного предельного цикла.
    \item \textbf{Прикладная математика и информатика} \textcolor{blue}{\textbf{[Запланировано]}}: Анализ бифуркаций вилки в нелинейных динамических системах первого и второго порядков.
    \item \textbf{Прикладная информатика (Мобильные технологии)} \textcolor{blue}{\textbf{[Запланировано]}}: Фазовые портреты динамики загрузки каналов связи в беспроводных Ad-Hoc сетях.
\end{itemize}

\subsection*{А.3. Лабораторная работа №3. Теория хаоса и функции Ляпунова}\label{app:A3}
\textit{Базовое задание:} Исследование нелинейных динамических систем, функций Ляпунова и чувствительности к начальным условиям.
\begin{itemize}[nosep]
    \item \textbf{Информационная безопасность} \textcolor{greenbordercolor}{\textbf{[Реализовано]}}: Моделирование аттрактора Лоренца $\dot{x} = \sigma(y-x), \; \dot{y} = x(\rho-z)-y, \; \dot{z} = xy-\beta z$ и системы реакция--диффузия $u_t = D u_{xx} + u - u^3$. Проверка асимптотической устойчивости через функцию Ляпунова $V(u) = \frac{D}{2}\int (u_x)^2 dx + \frac{1}{4}\int (1-u^2)^2 dx$. Генерация 256-битного ключа (SHA-256) и побитовое XOR-шифрование $C_i = P_i \oplus K_i$.
    \item \textbf{Мехатроника и робототехника} \textcolor{greenbordercolor}{\textbf{[Реализовано]}}: Нелинейный осциллятор Дуффинга $\ddot{x} + \delta \dot{x} + \alpha x + \beta x^3 = \gamma \cos(\omega t)$ в задачах виброзащиты манипуляторов. Исследование возникновения странного аттрактора при нелинейном резонансе.
    \item \textbf{Прикладная математика и информатика} \textcolor{blue}{\textbf{[Запланировано]}}: Вычисление наибольшего показателя Ляпунова $\lambda_{max} = \lim_{t\to\infty} \frac{1}{t}\ln\frac{\|\delta x(t)\|}{\|\delta x(0)\|}$ для системы Лоренца и оценка фрактальной размерности странного аттрактора.
    \item \textbf{Программная инженерия (Нейротехнологии)} \textcolor{blue}{\textbf{[Запланировано]}}: Нелинейная динамика рекуррентных нейронных сетей (Reservoir Computing / Neural ODE) $\dot{x} = -\frac{1}{\tau}x + W \tanh(x) + W_{in}u$. Исследование перехода от стационарных состояний к хаосу.
    \item \textbf{Бизнес-информатика} \textcolor{blue}{\textbf{[Запланировано]}}: Нелинейная модель детерминированного биржевого хаоса (модель финансового рынка Гудвина--Мандрыкина). Построение аттракторов колебаний цен акций.
    \item \textbf{Прикладная информатика (Мобильные технологии)} \textcolor{blue}{\textbf{[Запланировано]}}: Хаотические колебания задержек сетевых пакетов в протоколе TCP при алгоритмах управления перегрузкой (Congestion Control).
\end{itemize}

\subsection*{А.4. Лабораторная работа №4. Вариационное исчисление и оптимизация}\label{app:A4}
\textit{Базовое задание:} Формализация вариационных задач, сведение к уравнениям Эйлера--Лагранжа, дискретизация и численный поиск экстремалей.
\begin{itemize}[nosep]
    \item \textbf{Мехатроника и робототехника} \textcolor{greenbordercolor}{\textbf{[Реализовано]}}: Оптимизация траектории робота $y(x)$ на рельефе $z = h(x, y)$ путем минимизации функционала стоимости движения $\mathcal{J}[y] = \int_{a}^{b} (A + B |\nabla h|) \sqrt{1 + (y')^2} dx + C \int_{a}^{b} (y'')^2 dx \to \min$ методом градиентного спуска.
    \item \textbf{Информационная безопасность} \textcolor{greenbordercolor}{\textbf{[Реализовано]}}: Вариационная задача оптимального распределения мощности шифрования для минимизации риска утечки трафика $\mathcal{J}[P] = \int_0^T (\alpha Risk(t) + \beta P'(t)^2) dt \to \min$.
    \item \textbf{Бизнес-информатика} \textcolor{blue}{\textbf{[Запланировано]}}: Оптимальное управление производственным выпуском $u(t)$ и уровнем товарных запасов $x(t)$: $\mathcal{J}[u] = \int_0^T \left( c_1 (x(t) - x_{opt})^2 + c_2 u(t)^2 \right) dt \to \min$ при дифференциальной связи $\dot{x} = u(t) - d(t)$.
    \item \textbf{Прикладная математика и информатика} \textcolor{blue}{\textbf{[Запланировано]}}: Вариационная задача нахождения геодезических линий на римановых поверхностях $\mathcal{J}[y] = \int_a^b \sqrt{g_{11} + 2g_{12}y' + g_{22}(y')^2} dx \to \min$ и решение уравнения Эйлера--Лагранжа.
    \item \textbf{Программная инженерия (Нейротехнологии)} \textcolor{blue}{\textbf{[Запланировано]}}: Вариационный принцип обучения непрерывных Neural ODE: минимизация функционала потерь $\mathcal{J}[\theta] = L(x(T)) + \int_0^T R(x(t), \theta) dt$ с помощью сопряжённого метода (Adjoint State Method).
    \item \textbf{Прикладная информатика (Мобильные технологии)} \textcolor{blue}{\textbf{[Запланировано]}}: Оптимизация расхода энергии аккумуляторов мобильных устройств при непрерывной передачи данных $\mathcal{J}[P] = \int_0^T (P(t)^2 + \gamma \dot{P}(t)^2) dt \to \min$.
\end{itemize}

\newpage

\section*{Приложение Б. Покомпонентная система оценивания}\label{app:B}

\begin{greenblock}
\textbf{Накопительная система оценивания:} Практические элементы (РГР и лабораторные работы) составляют 60\% итоговой оценки каждого семестра.
\end{greenblock}

\vspace{0.5em}
\noindent\textbf{Таблица Б.1. Структура баллов за контрольные элементы (Семестр 1 и Семестр 2):}

\vspace{0.3em}
\begin{center}
\small
\begin{tabular}{|p{9.5cm}|c|c|}
\hline
\textbf{Элемент контроля} & \textbf{Баллы (Сем. 1)} & \textbf{Баллы (Сем. 2)} \\
\hline
РГР №1 / РГР №3 (профильно-адаптированная расчётная работа) & 10 & 10 \\
РГР №2 / РГР №4 (профильно-адаптированная расчётная работа) & 10 & 10 \\
Лабораторная работа №1 / Лабораторная работа №3 (профильный модуль) & 10 & 10 \\
Лабораторная работа №2 / Лабораторная работа №4 (профильный модуль) & 10 & 10 \\
Контрольная работа №1 (середина семестра) & 5 & 5 \\
Контрольная работа №2 (конец семестра) & 5 & 5 \\
Коллоквиум (теоретический минимум) & 20 & 20 \\
Итоговый экзамен & 20 & 20 \\
Работа на практических занятиях (активность, решение задач у доски) & 10 & 10 \\
\hline
\textbf{Итого} & \textbf{100} & \textbf{100} \\
\hline
\end{tabular}
\end{center}

\newpage

\section*{Список литературы и источников}\label{app:refs}

\subsection*{Фундаментальные учебники и монографии по ОДУ}
\begin{enumerate}[nosep, leftmargin=1.8em]
    \item \textbf{Степанов В. В.} Курс дифференциальных уравнений. — М.: Физматлит, 2008. — 512 с.
    \item \textbf{Филиппов А. Ф.} Введение в теорию дифференциальных уравнений. — М.: КомКнига, 2007. — 240 с.
    \item \textbf{Понтрягин Л. С.} Обыкновенные дифференциальные уравнения. — М.: Наука, 1982. — 331 с.
    \item \textbf{Арнольд В. И.} Обыкновенные дифференциальные уравнения. — М.: Наука, 1984. — 272 с.
\end{enumerate}

\subsection*{Устойчивость, теория хаоса и вариационное исчисление}
\begin{enumerate}[nosep, leftmargin=1.8em]
    \item \textbf{Кузнецов С. П.} Динамический хаос (лекционные курсы). — М.: Физматлит, 2006. — 356 с.
    \item \textbf{Хакен Г.} Синергетика: Иерархии неустойчивостей в самоорганизующихся системах. — М.: Мир, 1985. — 423 с.
    \item \textbf{Эльсгольц Л. Э.} Дифференциальные уравнения и вариационное исчисление. — М.: Наука, 1969. — 424 с.
    \item \textbf{Понтрягин Л. С., Болтянский В. Г., Гамкрелидзе Р. В., Мищенко Е. Ф.} Математическая теория оптимальных процессов. — М.: Наука, 1983. — 392 с.
\end{enumerate}

\subsection*{Вычислительные методы, ИИ и официальные ресурсы}
\begin{enumerate}[nosep, leftmargin=1.8em]
    \item \textbf{Chen R. T. Q., Rubanova Y., Bettencourt J., Duvenaud D. K.} Neural Ordinary Differential Equations // Advances in Neural Information Processing Systems (NeurIPS). — 2018. — Vol. 31. — \href{https://arxiv.org/abs/1806.07366}{arXiv:1806.07366}.
    \item \textbf{Официальный образовательный портал Университета ИТМО}: \href{https://abit.itmo.ru/}{https://abit.itmo.ru/} (Перечень образовательных программ бакалавриата).
\end{enumerate}

\end{document}